\begin{statementbox}
\textbf{条件与记号。}双曲线为
\[
 \frac{x^2}{a^2}-\frac{y^2}{b^2}=1\quad\text{或}\quad
 \frac{y^2}{a^2}-\frac{x^2}{b^2}=1,\qquad a,b>0.
\]
$a$ 为实半轴，$b$ 为虚半轴，$c=\sqrt{a^2+b^2}$，焦距为 $2c$。
$P$ 在曲线上且不是实轴顶点，$\theta=\angle F_1PF_2\in(0,\pi)$，
$S$ 为整个焦点三角形 $\triangle F_1PF_2$ 的面积。

\medskip
\textbf{核心结论。}
\[
 \boxed{S=b^2\cot\frac\theta2}
 =\boxed{\frac{b^2}{\tan(\theta/2)}}
 =\boxed{b^2\frac{\sin\theta}{1-\cos\theta}}.
\]
\textbf{辅助关系。}记 $m=PF_1$、$n=PF_2$，$h$ 为 $P$ 到焦点连线的垂距，则
\[
 |m-n|=2a,\qquad
 \boxed{b^2=mn\sin^2\frac\theta2},\qquad
 \boxed{S=ch}.
\]
横向标准式 $h=|y|$，纵向标准式 $h=|x|$；平移、旋转后仍须取到焦点连线的垂距。

\medskip
\textbf{范围与退化。}固定双曲线上，$0<\theta<\pi$ 均可实现，面积取遍 $(0,+\infty)$。
趋近顶点时 $\theta\to\pi$、$S\to0$；顶点本身位于两焦点之间，
向两焦点的射线反向，$\theta=\pi$，三角形退化，面积式可连续延伸为
$0=b^2\cot(\pi/2)$。

\medskip
\textbf{旁切圆几何提示（对应方法五）。}
把离 $P$ 较远的焦点重新记为 $F_1$，较近的记为 $F_2$，
使 $PF_1-PF_2=2a$。作三角形 $F_1PF_2$ 的
\textbf{$F_1$ 对边 $PF_2$ 所对应的旁切圆}，半径记为 $r$；
圆与线段 $PF_2$ 切于 $T$，则
\[
 \boxed{PT=c-a},\qquad
 \boxed{r=(c-a)\cot\frac\theta2},\qquad
 \boxed{S=(a+c)r}.
\]
三个小结论一乘：$S=(c+a)(c-a)\cot(\theta/2)=b^2\cot(\theta/2)$。
此法不使用余弦定理；原有计算方法仍可使用。

\medskip
\textbf{与椭圆区分。}
\[
 \text{椭圆： }S=b^2\tan\frac\theta2,
 \qquad \text{双曲线： }S=b^2\cot\frac\theta2.
\]
双曲线的 $a,b$ 无大小限制；面积式不显含 $a$，不等于焦点位置、点 $P$ 与实际夹角不受曲线约束。
\end{statementbox}
\clearpage
